\section{Background}
\TemplateTodo{Explain the general context, ideal condition, real problem, research/practical gap, and why the study is necessary.}

The number of smartphone users in Indonesia continues to grow rapidly each year. According to a We Are Social report, Indonesia had 212.9 million smartphone users in 2024 \cite{manning2008}. This growth has encouraged the development of mobile-based services, including online transportation and food delivery applications. GoFood, as one of the largest services in Indonesia, has millions of reviews on the Google Play Store that reflect user opinions and experiences.

User reviews contain valuable information for application developers to improve service quality. However, analyzing thousands of reviews manually requires significant time and effort. Therefore, an automated approach is needed to classify review sentiment into positive and negative categories.

Sentiment analysis is a branch of Natural Language Processing (NLP) that can be used to extract opinions from text automatically \cite{feldman2013}. The Na\"ive Bayes Classifier (NBC) method was chosen because it is simple, fast, and has proven effective in text classification \cite{go2019}. This study applies NBC with TF-IDF features to classify the sentiment of GoFood application reviews.
